\begin{tikzpicture}[
  font=\small,
  bx/.style={draw, rounded corners=1pt, inner sep=3pt, minimum height=5mm,
             minimum width=17mm, align=center, font=\small},
  ar/.style={-{Latex[length=1.7mm]}, line width=0.5pt},
  ttl/.style={font=\scriptsize\itshape, align=center},
  txt/.style={font=\scriptsize, align=center, inner sep=1.5pt}]

% =====================================================================
% (a)  the compactification:  P fixed, t in the closed simplex
% =====================================================================
\begin{scope}
\coordinate (v1) at (0,0);
\coordinate (v2) at (4.6,0);
\coordinate (v3) at (2.3,3.98);
\coordinate (w1) at (0.42,0.243);
\coordinate (w2) at (4.18,0.243);
\coordinate (w3) at (2.3,3.494);

\fill[black!7] (w1) -- (w2) -- (w3) -- cycle;
\fill[pattern=north east lines, pattern color=black!45, even odd rule]
  (v1) -- (v2) -- (v3) -- cycle   (w1) -- (w2) -- (w3) -- cycle;
\draw[line width=1pt] (v1) -- (v2) -- (v3) -- cycle;
\draw[black!35, line width=0.3pt] (w1) -- (w2) -- (w3) -- cycle;
\foreach \p in {(v1),(v2),(v3)}{\fill \p circle (1.5pt);}

\node[txt] at (2.3,2.25) {interior\\ $t_c>0$};
\fill (2.3,1.42) circle (1.7pt);
\node[txt, anchor=north] at (2.3,1.26) {$(P_{\min},t_{\min})$};
\node[txt, anchor=north] at (2.3,-0.06) {$t_z=0$};

\node[txt, anchor=north] at (2.3,-0.52)
  {closed simplex $\Delta^{m-1}$: $t_c\ge0$, $\sum_c t_c=1$;\\[1pt]
   the hatched faces $t_c=0$ carry no counterexample,\\[1pt]
   so a minimiser of $G$ lies in the interior};
\node[txt, anchor=north] at (2.3,-1.72)
  {$G(P,t)=\displaystyle\max_{|C|=k}\lmin\bigl(W[C]\bigr)$,\quad $W=P-\diag t$,\\[2pt]
   continuous on a compact set; the minimum is attained};
\end{scope}

% =====================================================================
% (b)  the two boundary mechanisms at a vanishing weight
% =====================================================================
\begin{scope}[xshift=5.9cm]
\node[txt, anchor=north west, text width=79mm, align=left] at (0,5.05)
  {A chart point at $(m,k)$ with $t_z=0$. The pivot identity
   $\lVert Pe_z\rVert^{2}=P_{zz}$ splits two cases.};

% ---- case A : the atom vanishes -------------------------------------
\node[ttl] at (1.95,4.05) {atom deletion};
\node[bx] (a1) at (1.95,3.62) {$P_{zz}=0$};
\node[txt, text width=38mm, align=center] (a2) at (1.95,2.52)
  {the column $Pe_z$ vanishes;\\ delete $z$, rescale $t\mapsto t/(1-t_z)$};
\node[bx] (a3) at (1.95,1.40) {$(m-1,k)$};
\node[txt, text width=38mm, align=center] (a4) at (1.95,0.78) {$C$ lifts verbatim};
\node[bx] (a5) at (1.95,0.15) {$|C|=k$};
\draw[ar] (a1) -- (a2); \draw[ar] (a2) -- (a3);
\draw[ar] (a3) -- (a4); \draw[ar] (a4) -- (a5);

% ---- case B : deflation at the pivot --------------------------------
\node[ttl] at (6.05,4.05) {Schur deflation};
\node[bx] (b1) at (6.05,3.62) {$P_{zz}>0$};
\node[txt, text width=38mm, align=center] (b2) at (6.05,2.52)
  {$P'=P-(Pe_z)(Pe_z)\tp/P_{zz}$,\\ $\tr P'=k-1$, $P'e_z=0$; delete $z$};
\node[bx] (b3) at (6.05,1.40) {$(m-1,k-1)$};
\node[txt, text width=38mm, align=center] (b4) at (6.05,0.78) {$C\mapsto C\cup\{z\}$};
\node[bx] (b5) at (6.05,0.15) {$|C\cup\{z\}|=k$};
\draw[ar] (b1) -- (b2); \draw[ar] (b2) -- (b3);
\draw[ar] (b3) -- (b4); \draw[ar] (b4) -- (b5);

% ---- the identity that pays for the second case ---------------------
\node[txt, anchor=north] at (4.0,-0.45)
  {$x\tp Wx=x\tp(P'-\diag t)\,x
    +\dfrac{\langle Pe_z,x\rangle^{2}}{P_{zz}}$,\\[3pt]
   and at $t_z=0$ the three terms complete a square};
\end{scope}

\end{tikzpicture}
